%! AP Statistics — Unit 3, Lesson 3.5 — Lesson Plan
\documentclass[10pt]{article}
\usepackage{apstats-article}
\usepackage{apstats-boxes}

\newcommand{\UnitNumberName}{Unit 3: Collecting Data \quad}
\newcommand{\LessonNumberName}{Lesson 3.5: Introduction to Experimental Design}

\begin{document}
\begin{center}
    {\Large\bfseries \CourseName: \SchoolYear\ (\MeetingLength) \\
    {\small\bfseries \UnitNumberName \LessonNumberName}}
\end{center}

\begin{tcolorbox}[colback=sky, colframe=navy, boxrule=0.9pt, arc=2mm,
    left=2mm, right=2mm, top=1.5mm, bottom=1.5mm]
    \textbf{Primary Objective:} Students will identify the components of a well-designed
    experiment, describe the four principles of experimental design, and compare
    completely randomized, block, and matched-pairs designs (Big Idea: VAR).
\end{tcolorbox}

\renewcommand{\arraystretch}{1.6}
\begin{skillbox}[Priority Ideas \& Skills]{goldbox}
\begin{minipage}[t]{0.48\linewidth}
    \textbf{AP Skill 1 --- Selecting Statistical Methods}
    \begin{itemize}
        \item Identify the components of an experiment (VAR-3.A, Skill 1.C).
        \item Describe elements of a well-designed experiment (VAR-3.B, Skill 1.B).
        \item Compare experimental designs: CRD, block, matched pairs (VAR-3.C,
              Skill 1.C).
    \end{itemize}
\end{minipage}
\hfill
\begin{minipage}[t]{0.48\linewidth}
    \textbf{Key Understandings}
    \begin{itemize}
        \item Four principles: comparison, random assignment, replication, control.
        \item Random \emph{assignment} (not selection) is what allows causal conclusions.
        \item Completely randomized design: all units randomly assigned to treatments.
        \item Block design: units grouped by a blocking variable; randomize within
              blocks to reduce variability.
        \item Matched pairs: a special block design with pairs of similar units.
    \end{itemize}
\end{minipage}
\end{skillbox}

\begin{skillbox}[{Vocabulary, Concepts, \& Theorems}]{greenbox}
\begin{minipage}[t]{0.48\linewidth}
    \begin{itemize}
        \item \textbf{Experimental unit} --- the individual assigned to a treatment
        \item \textbf{Factor (explanatory variable)} --- a variable whose levels are
              intentionally manipulated
        \item \textbf{Treatment} --- a specific level or combination of levels of the
              factor(s)
        \item \textbf{Response variable} --- the outcome measured after treatments
        \item \textbf{Confounding variable} --- related to factor and response; can
              create a false perception of association
        \item \textbf{Control group} --- receives no treatment or a placebo
        \item \textbf{Placebo} --- an inactive treatment; controls for the placebo effect
    \end{itemize}
\end{minipage}
\hfill
\begin{minipage}[t]{0.48\linewidth}
    \begin{itemize}
        \item \textbf{Random assignment} --- experimental units assigned to treatments
              by chance; balances confounding variables
        \item \textbf{Replication} --- more than one experimental unit per treatment group
        \item \textbf{Single-blind} --- subjects don't know which treatment they receive
        \item \textbf{Double-blind} --- neither subjects nor evaluators know which
              treatment subjects receive
        \item \textbf{Block design} --- experimental units grouped by a blocking
              variable; treatments assigned randomly within each block
        \item \textbf{Matched pairs} --- a block design with two units per block (or
              same unit receives both treatments)
    \end{itemize}
\end{minipage}
\end{skillbox}

\begin{skillbox}[Activate Prior Knowledge \& Spiral Review]{sky}
\begin{minipage}[t]{0.48\linewidth}
    \textbf{Prior Knowledge Check (3--4 min)}
    \begin{itemize}
        \item Recall: ``What is the difference between an observational study and an
              experiment?'' (Lesson 3.2)
        \item Recall: ``What is a confounding variable?'' (Unit 2, Lesson 3.2)
        \item Transition: ``Experiments are designed to \emph{eliminate} confounding.
              Today we learn how.''
    \end{itemize}
\end{minipage}
\hfill
\begin{minipage}[t]{0.48\linewidth}
    \textbf{Connections Forward}
    \begin{itemize}
        \item Choosing the right design for a given situation $\to$ Lesson 3.6
        \item Interpreting experiment results; statistical significance $\to$ Lesson 3.7
        \item Significance testing for experiments $\to$ Units 6--9
    \end{itemize}
\end{minipage}
\end{skillbox}

\begin{skillbox}[Hook \& Lesson]{sky}
\begin{minipage}[t]{0.48\linewidth}
    \textbf{Hook (5--7 min)}\\
    Scenario: \textit{``I want to test whether a new study app improves test scores.
    I let students choose whether to use it. App users score 8 points higher. Can I
    conclude the app works?''}\\[4pt]
    Students discuss: ``What's wrong with this study?'' (Selection bias: motivated
    students choose to use the app.)\\[4pt]
    Transition: \textit{``An experiment randomly assigns which students use the app,
    eliminating that problem.''}
\end{minipage}
\hfill
\begin{minipage}[t]{0.48\linewidth}
    \textbf{Lesson Flow (15--18 min)}
    \begin{enumerate}
        \item Identify components: experimental units, factor, treatments, response,
              confounding variable.
        \item Four principles of well-designed experiments (VAR-3.B.1):
              (a) comparison, (b) random assignment, (c) replication,
              (d) control of confounding.
        \item Completely randomized design: all units randomly assigned to treatments;
              diagram.
        \item Blinding: single-blind and double-blind; placebo effect.
        \item Block design: reduce variability by blocking on a variable (e.g., gender);
              randomize within blocks.
        \item Matched pairs: pairs of similar units; one per treatment. OR same unit
              gets both treatments in random order.
    \end{enumerate}
\end{minipage}
\end{skillbox}

\begin{skillbox}[Explicit Instruction \& Active Monitoring]{sky}
\begin{minipage}[t]{0.48\linewidth}
    \textbf{Direct Instruction (10 min)}
    \begin{itemize}
        \item Worked example: testing a new fertilizer on corn plants. Identify
              units, factor, treatments, response, and potential confounds. Design a
              CRD, then a block design blocking on soil type.
        \item Key distinction: blocking variable is a known source of variability that
              we want to \emph{remove} from the error, not a treatment we are testing.
    \end{itemize}
\end{minipage}
\hfill
\begin{minipage}[t]{0.48\linewidth}
    \textbf{Active Monitoring}
    \begin{itemize}
        \item Listen for: confusing random \emph{assignment} with random \emph{selection}.
        \item Cold-call: ``If I don't randomly assign, which principle of good
              experimental design am I violating?''
        \item Check that students can draw a diagram of the design (units $\to$ random
              assignment $\to$ treatment groups $\to$ response).
    \end{itemize}
\end{minipage}
\end{skillbox}

\begin{skillbox}[Group Work \& Differentiation]{redbox}
\begin{minipage}[t]{0.48\linewidth}
    \textbf{Partner Activity (8--10 min)}\\
    Three experiment descriptions. Partners:
    \begin{enumerate}
        \item Identify experimental units, factor(s), treatments, and response.
        \item Identify the design type (CRD, block, or matched pairs).
        \item Identify whether blinding is used and what type.
        \item State one potential confounding variable and how the design controls it.
    \end{enumerate}
\end{minipage}
\hfill
\begin{minipage}[t]{0.48\linewidth}
    \textbf{Differentiation}
    \begin{itemize}
        \item \textit{Support}: Diagram template --- boxes for units, random
              assignment arrow, treatment boxes, response measurement.
        \item \textit{Extension}: Design a double-blind experiment to test whether a
              new energy drink improves reaction time. Explain each design choice.
        \item \textit{ELL}: Visual experiment diagram with labels in both English and
              L1 if available.
    \end{itemize}
\end{minipage}
\end{skillbox}

\begin{skillbox}[Individual Work \& Assessment]{redbox}
\begin{minipage}[t]{0.48\linewidth}
    \textbf{Exit Ticket (5 min)}\\
    Scenario: \textit{``Researchers recruit 60 volunteers to test a new headache
    medication. Half receive the drug; half receive a sugar pill. Neither group
    knows which they received. Researchers measure headache severity after 1 hour.''}
    \begin{enumerate}
        \item Identify: experimental units, treatment(s), response variable.
        \item What design is this? (CRD, block, or matched pairs?)
        \item What type of blinding is used?
        \item Why is random assignment essential?
    \end{enumerate}
\end{minipage}
\hfill
\begin{minipage}[t]{0.48\linewidth}
    \textbf{AP-Style Check}\\
    \textit{``In an experiment testing whether music improves focus, students are
    randomly assigned to study with music or in silence for 30 minutes, then take
    a quiz. Researchers block by grade level (9th vs.\ 12th). What is the purpose
    of blocking?''}
    \begin{enumerate}[label=(\Alph*)]
        \item To create larger treatment groups
        \item To reduce variability due to grade-level differences
        \item To ensure all students take the same quiz
        \item To prevent the placebo effect
    \end{enumerate}
    Answer: (B)
\end{minipage}
\end{skillbox}

\begin{skillbox}[Reinforcement \& Extension]{goldbox}
\begin{minipage}[t]{0.48\linewidth}
    \textbf{Homework / Practice}
    \begin{itemize}
        \item For 4 experiment descriptions: identify units/factor/treatments/response,
              classify the design, identify type of blinding, and evaluate whether the
              four principles of good design are met.
        \item Reflection: \textit{``Why does random \emph{assignment} (not just random
              selection) allow us to conclude causation?''}
    \end{itemize}
\end{minipage}
\hfill
\begin{minipage}[t]{0.48\linewidth}
    \textbf{Extension / Preview}
    \begin{itemize}
        \item \textit{Extension}: Research a famous clinical trial (e.g., Salk
              polio vaccine, 1954). Identify the components of the experimental design
              and explain why it was double-blind.
        \item \textit{Preview}: Lesson 3.6 asks you to choose the \emph{best} design
              for a given situation --- think about what factors matter when choosing
              between CRD, block, and matched pairs.
    \end{itemize}
\end{minipage}
\end{skillbox}

\end{document}
